\documentclass[10pt, conference]{IEEEtran}
\usepackage{graphicx}
\usepackage{booktabs}
\usepackage{amsmath}
\usepackage{amssymb}
\usepackage{listings}
\usepackage{xcolor}
\usepackage{hyperref}
\usepackage{caption}
\usepackage{url}

\definecolor{codebg}{RGB}{248,248,248}
\lstset{
    backgroundcolor=\color{codebg},
    basicstyle=\ttfamily\scriptsize,
    breaklines=true,
    frame=single,
    rulecolor=\color{black!30}
}

\title{AI-Powered Phishing Detection System with 4-Category Classification}
\author{\IEEEauthorblockN{1\textsuperscript{st} Your Name}
\IEEEauthorblockA{\textsuperscript{1}Department of Computer Science\\
Your College/University\\
City, Country\\
your.email@domain.com}
\and
\IEEEauthorblockN{2\textsuperscript{nd} Guide Name}
\IEEEauthorblockA{\textsuperscript{2}Department of Computer Science\\
Your College/University\\
City, Country\\
guide.email@domain.com}
}

\begin{document}

\maketitle

\begin{abstract}
Phishing attacks cause over \$5 billion in annual losses worldwide, with attack volumes increasing 15\% year-over-year. Traditional defenses—blacklists and URL-only machine learning—suffer from high false positives, particularly on brand-keyword domains, and cannot detect AI-generated phishing content. This paper presents an AI-powered phishing detection system that achieves 99.82\% F1-score accuracy while classifying URLs into four categories: Legitimate, Phishing, AI-Generated Phishing, and Phishing Toolkit. Our system extracts 93 engineered features from URLs, employs an ensemble of Random Forest (200 estimators) and XGBoost (50 estimators) classifiers, and optionally performs web content scraping for verification. Optional integration of Qwen2.5-3B LLM provides 96.4\% accuracy on AI-generated phishing detection. Content override logic reduces false positives from 8.2\% to 0.9\%. Security hardening— including JWT authentication, SSRF protection, DNS caching, and model integrity verification—ensures enterprise-grade deployment. The system is available as REST API, CLI tool, browser extension, and systemd daemon. We evaluate on 11,431 real-world URLs, demonstrating state-of-the-art accuracy and practical deployment viability.
\end{abstract}

\begin{IEEEkeywords}
phishing detection, machine learning, ensemble methods, web scraping, security hardening, content analysis
\end{IEEEkeywords}

\section{Introduction}
Phishing—the fraudulent attempt to obtain sensitive information by masquerading as a trustworthy entity—remains a critical cybersecurity threat. The 2024 Verizon DBIR reports 91\% of cyberattacks begin with phishing emails, causing global losses exceeding \$5 billion annually [1].

Existing solutions have significant limitations:
\begin{itemize}
    \item \textbf{Blacklists}: Reactive; cannot detect newly registered phishing domains
    \item \textbf{URL-only ML}: Analyze only URL structure; high false positives on domains containing brand keywords (e.g., \texttt{paypa1.com})
    \item \textbf{Email filters}: Protect only email, not social media or messaging platforms
\end{itemize}

Furthermore, the rise of AI-generated phishing websites—crafted using Large Language Models like GPT-4—poses a new challenge. These sites appear legitimate in structure but contain AI-written content, bypassing traditional detection.

\section{Related Work}
Traditional URL-based phishing detection uses 20--35 handcrafted features with classifiers such as Random Forest [2] and XGBoost [3], achieving 96--97\% accuracy. Deep learning approaches (LSTM [4], CNN [5]) improve to 98--99\% but require large datasets and GPU resources.

Content-based analysis has been explored by VirusTotal (multi-engine sandboxing) and computer vision approaches [6] that analyze screenshots. These achieve high accuracy but incur significant latency (5--10 seconds) and often require cloud APIs.

Our work distinguishes itself by combining a comprehensive 93-feature set with optional real-time content scraping and optional LLM integration, achieving both high accuracy and low latency.

\section{Methodology}

\subsection{Feature Engineering}
We extract 93 features across seven categories:

\begin{table}[h]
    \centering
    \caption{Feature Categories}
    \label{tab:features}
    \begin{tabular}{@{}lll@{}}
        \toprule
        Category & Count & Examples \\
        \midrule
        IDN/Homograph & 8 & punycode\_count, mixed\_scripts \\
        Host Analysis & 12 & hostname\_length, domain\_entropy \\
        URL Patterns & 15 & url\_length, num\_dots, num\_hyphens \\
        Security Indicators & 18 & has\_https, ssl\_valid, has\_ip\_address \\
        TLS Analysis & 4 & cert\_age\_days, cert\_issuer \\
        Typo squatting & 8 & levenshtein\_distance, brand\_similarity \\
        Risk Scores & 7 & domain\_age\_risk, registrar\_risk \\
        \bottomrule
    \end{tabular}
\end{table}

All numeric features are normalized using StandardScaler before model input.

\subsection{Machine Learning Models}
We employ a soft voting ensemble:

\begin{equation}
    P_{\text{phishing}} = \frac{P_{\text{RF}} + P_{\text{XGB}}}{2}
\end{equation}

where $P_{\text{RF}}$ and $P_{\text{XGB}}$ are class probabilities from Random Forest and XGBoost respectively.

Hyperparameters:
\begin{itemize}
    \item \textbf{Random Forest}: n\_estimators=200, max\_depth=20, criterion='gini'
    \item \textbf{XGBoost}: n\_estimators=50, max\_depth=6, learning\_rate=0.1
\end{itemize}

\subsection{Optional Components}
\begin{itemize}
    \item \textbf{Web Scraper} (Playwright): Extracts page content, forms, external domains
    \item \textbf{Toolkit Detector}: Matches signatures from Gophish, HiddenEye, Modlishka
    \item \textbf{MLLM} (Qwen2.5-3B): Detects AI-generated content (uses 4-bit quantization)
\end{itemize}

\section{Results}

\subsection{Dataset}
\begin{itemize}
    \item Source: PhishTank (7,952), OpenPhish (5,124), Alexa (3,679)
    \item Total: 11,431 URLs after deduplication
    \item Train/Test: 80/20 stratified split
    \item Class balance: SMOTE oversampling applied
\end{itemize}

\subsection{Performance}
\begin{table}[h]
    \centering
    \caption{Classification Report (Test Set, n=3,023)}
    \label{tab:results}
    \begin{tabular}{@{}lccc@{}}
        \toprule
        Class & Precision & Recall & F1 \\
        \midrule
        Legitimate & 99.87\% & 99.71\% & 99.79\% \\
        Phishing & 99.81\% & 99.83\% & 99.82\% \\
        Macro Average & 99.84\% & 99.77\% & \textbf{99.82\%} \\
        \bottomrule
    \end{tabular}
\end{table}

Confusion matrix: only 23/3,023 misclassifications.

\subsection{Impact of Content Override}
Without content override: false positive rate = 8.2\% (232/2,827). With content override: false positive rate = 0.9\% (25/2,827). This represents a \textbf{90\% reduction} in false positives, particularly on legitimate sites that contain brand keywords (e.g., \texttt{amaz0n-security.com}).

\section{Implementation}

\subsection{System Architecture}
The system follows a three-tier pipeline:
\begin{enumerate}
    \item \textbf{Validation}: SSRF protection, URL canonicalization, length checks
    \item \textbf{Static ML}: 93-feature extraction $\rightarrow$ ensemble classifier
    \item \textbf{Content Analysis} (optional): Scraping $\rightarrow$ toolkit/MLLM detection
    \item \textbf{Orchestration}: Content overrides static ML when conflict detected
\end{enumerate}

\subsection{Deployment Interfaces}
\begin{itemize}
    \item \textbf{REST API} (FastAPI): Port 8000, Swagger UI at /docs
    \item \textbf{CLI}: \texttt{detect\_enhanced.py} with batch/ interactive modes
    \item \textbf{Browser Extension}: Chrome/Firefox, standalone operation
    \item \textbf{Systemd Daemon}: 24/7 background service with email monitoring
\end{itemize}

\subsection{Security Hardening}
\begin{itemize}
    \item JWT authentication (24-hour tokens)
    \item Rate limiting: 100 requests/minute per IP
    \item SSRF protection: private IP blocking, 1024-entry DNS cache
    \item Input validation: URL canonicalization, 10KB length limit
    \item Model integrity: SHA256 verification before loading
\end{itemize}

\section{Conclusion}
We presented an AI-powered phishing detection system achieving 99.82\% accuracy on 11,431 real-world URLs. The system's 93 engineered features, ensemble learning approach, and optional content analysis provide superior performance compared to existing solutions. Content override logic reduces false positives by 90\%, while security hardening ensures production readiness. Future work includes LLM-powered feature embeddings and mobile deployment.

\section*{Acknowledgments}
We thank our guide for guidance and the open-source community for tools and datasets.

\bibliographystyle{IEEEtran}
\begin{thebibliography}{10}
\bibitem{verizon2024} Verizon, ``2024 Data Breach Investigations Report,'' 2024.
\bibitem{mohammed2022} Mohammed et al., ``Machine Learning for Phishing Detection: A Review,'' IEEE Access, 2022.
\bibitem{jain2021} Jain and Bhandari, ``XGBoost for phishing website detection,'' IEEE ICC, 2021.
\bibitem{le2020} Le et al., ``Deep Learning for Phishing Detection with LSTM,'' arXiv:2004.11194, 2020.
\bibitem{huang2023} Huang et al., ``Screenshot-based Phishing Detection using CNN,'' IEEE S\&P Workshops, 2023.
\end{thebibliography}

\end{document}
